% Anhang
\appendix
\section{User Research}
\label{app:user-research}

\subsection{Interview-Leitfaden (Auszug)}

Im Zeitraum KW~18--19 führte die Gruppe zehn semi-strukturierte Interviews (N=10) mit Studierenden des Studiengangs Angewandte Informatik und Soziale Medien (3.--6.~Semester) durch. Die Auswertung erfolgte anonymisiert ohne Nennung personenbezogener Namen. Der Leitfaden umfasste unter anderem:

\begin{enumerate}
    \item Was weißt du über das Immersive Reality Lab?
    \item Wenn du die Website besuchen würdest, was würdest du als erstes suchen?
    \item Was würde dich beim Suchen auf der Website frustrieren?
    \item Würdest du lieber suchen/klicken oder eine Frage stellen?
    \item Welche Art von Fragen würdest du einem Chatbot zum Lab stellen?
\end{enumerate}

\subsection{Personas}

\begin{tabularx}{\textwidth}{@{}l X X@{}}
\toprule
\textbf{Persona} & \textbf{Ziele} & \textbf{Fragen} \\
\midrule
Stud., 22 J., w & Abschlussarbeit, Ausleihe & VR-Projekte, Kurse \\
Forsch., 35 J., m & Publikationen, Kooperation & Biosensorik, Team \\
\bottomrule
\end{tabularx}

\subsection{Top-Informationsbedürfnisse}

{\tablecompact
\begin{tabularx}{\textwidth}{@{}c X r@{}}
\toprule
\textbf{\#} & \textbf{Bedürfnis} & \textbf{N} \\
\midrule
1 & Projekte im Lab & 10/10 \\
2 & Geräteausleihe & 10/10 \\
3 & Publikationen zu einem Thema & 8/10 \\
4 & Teammitglieder & 8/10 \\
5 & Abschlussarbeitsthemen & 6/10 \\
6 & Lehrveranstaltungen & 6/10 \\
7 & Kontakt zum Lab & 6/10 \\
8 & Technologien (VR/AR/Biosensorik) & 4/10 \\
9 & Methoden und Infrastruktur & 4/10 \\
\bottomrule
\end{tabularx}
}

\section{SUS-Rohdaten}
\label{app:sus}

{\scriptsize
\begin{tabularx}{\textwidth}{@{}X cccccccccc@{}}
\toprule
\textbf{Item} & \textbf{P1} & \textbf{P2} & \textbf{P3} & \textbf{P4} & \textbf{P5} & \textbf{P6} & \textbf{P7} & \textbf{P8} & \textbf{P9} & \textbf{P10} \\
\midrule
1. Häufige Nutzung & 3 & 5 & 3 & 3 & 4 & 3 & 3 & 3 & 3 & 4 \\
2. Unnötig komplex & 3 & 3 & 2 & 3 & 1 & 3 & 3 & 1 & 2 & 1 \\
3. Einfach zu nutzen & 4 & 3 & 4 & 3 & 5 & 5 & 3 & 5 & 4 & 4 \\
4. Brauche Unterstützung & 3 & 1 & 1 & 3 & 3 & 2 & 2 & 3 & 2 & 2 \\
5. Gut integriert & 4 & 3 & 5 & 4 & 5 & 5 & 5 & 3 & 4 & 5 \\
6. Inkonsistenzen & 2 & 2 & 2 & 2 & 1 & 1 & 2 & 1 & 3 & 2 \\
7. Schnell erlernbar & 5 & 3 & 4 & 5 & 5 & 5 & 3 & 4 & 4 & 3 \\
8. Umständlich & 2 & 3 & 3 & 3 & 2 & 3 & 1 & 1 & 1 & 3 \\
9. Sicher bei Nutzung & 5 & 3 & 4 & 5 & 4 & 3 & 3 & 4 & 5 & 4 \\
10. Viel vorher lernen & 2 & 2 & 1 & 1 & 3 & 2 & 1 & 3 & 3 & 2 \\
\midrule
\textbf{SUS} & \textbf{72,5} & \textbf{65,0} & \textbf{77,5} & \textbf{70,0} & \textbf{82,5} & \textbf{75,0} & \textbf{70,0} & \textbf{75,0} & \textbf{72,5} & \textbf{75,0} \\
\bottomrule
\end{tabularx}
}

\noindent Mittelwert: 73,5 (SD: 4,7; N=10). Berechnung nach Lewis~\cite{lewis2018}; Einzelwerte aus den Item-Antworten (Skala 1--5) nach Brooke.
